\documentclass[UTF8,12pt,fontset=windows]{ctexart}
\usepackage[paperwidth=240mm,paperheight=200mm,left=14mm,right=14mm,top=9mm,bottom=12mm]{geometry}
\usepackage{amsmath,amssymb,mathtools,bm,graphicx,xcolor}
\pagestyle{empty}
\setlength{\parindent}{0pt}
\setlength{\parskip}{4pt}
\setlength{\emergencystretch}{2em}
\xeCJKDeclareCharClass{CJK}{"2460 -> "24FF}
\xeCJKDeclareCharClass{CJK}{"2160 -> "2188}
\begin{document}
{\large\bfseries 2012.1}

曲线 $y = \frac{x^2 + x}{x^2 - 1}$ 的渐近线的条数为（　）

（A）0

（B）1

（C）2

（D）3

\vfill
\newpage
{\large\bfseries 2012.2}

设函数 $f(x) = (e^x - 1)(e^{2x} - 2)\cdots(e^{nx} - n)$，其中 $n$ 为正整数，则 $f'(0) =$（　）

（A）$(-1)^{n-1}(n-1)!$

（B）$(-1)^n(n-1)!$

（C）$(-1)^{n-1}n!$

（D）$(-1)^nn!$

\vfill
\newpage
{\large\bfseries 2012.3}

设 $a_n > 0$（$n = 1, 2, \cdots$），$S_n = a_1 + a_2 + \cdots + a_n$，则数列 $\{S_n\}$ 有界是数列 $\{a_n\}$ 收敛的（　）

（A）充分必要条件

（B）充分非必要条件

（C）必要非充分条件

（D）既非充分也非必要条件

\vfill
\newpage
{\large\bfseries 2012.4}

设 $I_k = \int_0^{k\pi} e^{x^2}\sin x dx$（$k = 1, 2, 3$），则有（　）

（A）$I_1 < I_2 < I_3$

（B）$I_3 < I_2 < I_1$

（C）$I_2 < I_3 < I_1$

（D）$I_2 < I_1 < I_3$

\vfill
\newpage
{\large\bfseries 2012.5}

设函数 $f(x, y)$ 可微，且对任意的 $x, y$ 都有 $\frac{\partial f(x, y)}{\partial x} > 0, \frac{\partial f(x, y)}{\partial y} < 0$，则使不等式 $f(x_1, y_1) < f(x_2, y_2)$ 成立的一个充分条件是（　）

（A）$x_1 > x_2, y_1 < y_2$

（B）$x_1 > x_2, y_1 > y_2$

（C）$x_1 < x_2, y_1 < y_2$

（D）$x_1 < x_2, y_1 > y_2$

\vfill
\newpage
{\large\bfseries 2012.6}

设区域 $D$ 由曲线 $y = \sin x, x = \pm\frac{\pi}{2}, y = 1$ 围成，则 $\iint_D (xy^5 - 1)dxdy =$（　）

（A）$\pi$

（B）2

（C）$-2$

（D）$-\pi$

\vfill
\newpage
{\large\bfseries 2012.7}

设 $\boldsymbol{\alpha}_1 = \begin{pmatrix} 0 \\ 0 \\ c_1 \end{pmatrix}, \boldsymbol{\alpha}_2 = \begin{pmatrix} 0 \\ 1 \\ c_2 \end{pmatrix}, \boldsymbol{\alpha}_3 = \begin{pmatrix} 1 \\ -1 \\ c_3 \end{pmatrix}, \boldsymbol{\alpha}_4 = \begin{pmatrix} -1 \\ 1 \\ c_4 \end{pmatrix}$，其中 $c_1, c_2, c_3, c_4$ 为任意常数，则下列向量组线性相关的为（　）

（A）$\boldsymbol{\alpha}_1, \boldsymbol{\alpha}_2, \boldsymbol{\alpha}_3$

（B）$\boldsymbol{\alpha}_1, \boldsymbol{\alpha}_2, \boldsymbol{\alpha}_4$

（C）$\boldsymbol{\alpha}_1, \boldsymbol{\alpha}_3, \boldsymbol{\alpha}_4$

（D）$\boldsymbol{\alpha}_2, \boldsymbol{\alpha}_3, \boldsymbol{\alpha}_4$

\vfill
\newpage
{\large\bfseries 2012.8}

设 $\boldsymbol{A}$ 为3阶矩阵，$\boldsymbol{P}$ 为3阶可逆矩阵，且 $\boldsymbol{P}^{-1}\boldsymbol{AP} = \begin{pmatrix} 1 & 0 & 0 \\ 0 & 1 & 0 \\ 0 & 0 & 2 \end{pmatrix}$。若 $\boldsymbol{P} = (\boldsymbol{\alpha}_1, \boldsymbol{\alpha}_2, \boldsymbol{\alpha}_3)$，$\boldsymbol{Q} = (\boldsymbol{\alpha}_1 + \boldsymbol{\alpha}_2, \boldsymbol{\alpha}_2, \boldsymbol{\alpha}_3)$，则 $\boldsymbol{Q}^{-1}\boldsymbol{AQ} =$（　）

（A）$\begin{pmatrix} 1 & 0 & 0 \\ 0 & 2 & 0 \\ 0 & 0 & 1 \end{pmatrix}$

（B）$\begin{pmatrix} 1 & 0 & 0 \\ 0 & 1 & 0 \\ 0 & 0 & 2 \end{pmatrix}$

（C）$\begin{pmatrix} 2 & 0 & 0 \\ 0 & 1 & 0 \\ 0 & 0 & 2 \end{pmatrix}$

（D）$\begin{pmatrix} 2 & 0 & 0 \\ 0 & 2 & 0 \\ 0 & 0 & 1 \end{pmatrix}$

\vfill
\newpage
{\large\bfseries 2012.9}

设 $y = y(x)$ 是由方程 $x^2 - y + 1 = e^y$ 所确定的隐函数，则 $\left.\frac{d^2y}{dx^2}\right|_{x = 0} =$ \underline{\hspace{20mm}}.

\vfill
\newpage
{\large\bfseries 2012.10}

$\lim_{n \to \infty}n\left(\frac{1}{1 + n^2} + \frac{1}{2^2 + n^2} + \cdots + \frac{1}{n^2 + n^2}\right) =$ \underline{\hspace{20mm}}.

\vfill
\newpage
{\large\bfseries 2012.11}

设 $z = f\left(\ln x + \frac{1}{y}\right)$，其中函数 $f(u)$ 可微，则 $x\frac{\partial z}{\partial x} + y^2\frac{\partial z}{\partial y} =$ \underline{\hspace{20mm}}.

\vfill
\newpage
{\large\bfseries 2012.12}

微分方程 $y dx + (x - 3y^2)dy = 0$ 满足条件 $y|_{x = 1} = 1$ 的解为 $y =$ \underline{\hspace{20mm}}.

\vfill
\newpage
{\large\bfseries 2012.13}

曲线 $y = x^2 + x$（$x < 0$）上曲率为 $\frac{\sqrt{2}}{2}$ 的点的坐标是 \underline{\hspace{20mm}}.

\vfill
\newpage
{\large\bfseries 2012.14}

设 $\boldsymbol{A}$ 为3阶矩阵，$|\boldsymbol{A}| = 3$，$\boldsymbol{A}^*$ 为 $\boldsymbol{A}$ 的伴随矩阵，若交换 $\boldsymbol{A}$ 的第1行与第2行得矩阵 $\boldsymbol{B}$，则 $|\boldsymbol{BA}^*| =$ \underline{\hspace{20mm}}.

\vfill
\newpage
{\large\bfseries 2012.15}

（本题满分10分）

已知函数 $f(x) = \frac{1 + x}{\sin x} - \frac{1}{x}$，记 $a = \lim_{x \to 0}f(x)$.

（Ⅰ）求 $a$ 的值；

（Ⅱ）若当 $x \to 0$ 时，$f(x) - a$ 与 $x^k$ 是同阶无穷小量，求常数 $k$ 的值.

\vfill
\newpage
{\large\bfseries 2012.16}

（本题满分10分）

求函数 $f(x, y) = xe^{-\frac{x^2 + y^2}{2}}$ 的极值.

\vfill
\newpage
{\large\bfseries 2012.17}

（本题满分12分）

过点 $(0, 1)$ 作曲线 $L: y = \ln x$ 的切线，切点为 $A$，又 $L$ 与 $x$ 轴交于 $B$ 点，区域 $D$ 由 $L$ 与直线 $AB$ 围成。求区域 $D$ 的面积及 $D$ 绕 $x$ 轴旋转一周所得旋转体的体积。

\vfill
\newpage
{\large\bfseries 2012.18}

（本题满分10分）

计算二重积分 $\iint_D xy d\sigma$，其中区域 $D$ 由曲线 $r = 1 + \cos\theta$（$0 \le \theta \le \pi$）与极轴围成。

\vfill
\newpage
{\large\bfseries 2012.19}

（本题满分10分）

已知函数 $f(x)$ 满足方程 $f'''(x) + f'(x) - 2f(x) = 0$ 及 $f''(x) + f(x) = 2e^x$.

（Ⅰ）求 $f(x)$ 的表达式；

（Ⅱ）求曲线 $y = f(x^2)\int_0^x f(-t^2)dt$ 的拐点。

\vfill
\newpage
{\large\bfseries 2012.20}

（本题满分10分）

证明：$x\ln\frac{1 + x}{1 - x} + \cos x \ge 1 + \frac{x^2}{2}$（$-1 < x < 1$）.

\vfill
\newpage
{\large\bfseries 2012.21}

（本题满分10分）

（Ⅰ）证明方程 $x^n + x^{n-1} + \cdots + x = 1$（$n$ 为大于1的整数）在区间 $\left(\frac{1}{2}, 1\right)$ 内有且仅有一个实根；

（Ⅱ）记（Ⅰ）中的实根为 $x_n$，证明 $\lim_{n \to \infty}x_n$ 存在，并求此极限.

\vfill
\newpage
{\large\bfseries 2012.22}

（本题满分11分）

设 $\boldsymbol{A} = \begin{pmatrix} 1 & a & 0 & 0 \\ 0 & 1 & a & 0 \\ 0 & 0 & 1 & a \\ a & 0 & 0 & 1 \end{pmatrix}, \boldsymbol{\beta} = \begin{pmatrix} 1 \\ -1 \\ 0 \\ 0 \end{pmatrix}$.

（Ⅰ）计算行列式 $|\boldsymbol{A}|$；

（Ⅱ）当实数 $a$ 为何值时，方程组 $\boldsymbol{Ax} = \boldsymbol{\beta}$ 有无穷多解，并求其通解.

\vfill
\newpage
{\large\bfseries 2012.23}

（本题满分11分）

已知 $\boldsymbol{A} = \begin{pmatrix} 1 & 0 & 1 \\ 0 & 1 & 1 \\ -1 & 0 & a \\ 0 & a & -1 \end{pmatrix}$，二次型 $f(x_1, x_2, x_3) = \boldsymbol{x}^T(\boldsymbol{A}^T\boldsymbol{A})\boldsymbol{x}$ 的秩为2.

（Ⅰ）求实数 $a$ 的值；

（Ⅱ）求正交变换 $\boldsymbol{x} = \boldsymbol{Qy}$ 将 $f$ 化为标准形.

\vfill
\end{document}
